\chapter{Where assistance earns its cost}
\label{reading:assistance-conclusions}

\section{Improve the question before automating the answer}

The strongest proposal for language assistance begins with tasks that a bank
economist can inspect: separating an invitation from actual participation,
identifying which information was available before a decision, proposing
alternative causal diagrams, and finding the source of an assumption. These
tasks can improve a study before any treatment effect is estimated. They also
provide concrete ways to compare assistance with a structured questionnaire,
a conventional search, or an unaided analyst.~\citep{v12_la_2305_00050,v12_la_2606_23608}

Text can supply useful information about mechanisms, but a plausible sentence
does not establish what happened in a bank. A process record showing that an
invitation was randomized has a different evidential role from a model's
general statement that banks sometimes randomize offers. Retrieval helps only
when the retrieved material is relevant and reliable. The reported LACR
comparisons include adverse as well as favourable retrieval outcomes, which
makes indiscriminate retrieval an inadequate design recommendation.
\citep{v12_la_2402_15301}

\section{Better measurement can still leave confounding}

Extracting financial distress from customer notes illustrates the distinction.
A model may summarize the notes accurately while the notes omit circumstances
that influenced both a servicing action and repayment. The measurement example
in Section~\ref{exec13:census-measurement} makes the consequence explicit.
With true distress $U$, measured distress $W=U+\varepsilon_W$, action
$A=U+\varepsilon_A$ and outcome $Y=\tau A+U+\varepsilon_Y$, mutually
independent centred Gaussian disturbances yield
\[
 \beta_{A\mid W}=\tau+\frac{s}{1+2s},\qquad
 s=\operatorname{Var}(\varepsilon_W),
\]
when the other variances are one. The second term is residual confounding.
It persists in an arbitrarily large sample. A narrower statistical interval
around the adjusted coefficient can therefore make an incorrect economic
interpretation appear more certain. The primary review of text as confounder
measurement discusses the need to distinguish confounders, mediators and
colliders; the displayed calculation is the manuscript's constructed example.
\citep{exec13_census_text_review_arxiv}

\begin{figure}[htbp]\centering
\begin{tikzpicture}[x=2.2cm,y=6cm,>=Stealth,every node/.style={font=\small}]
 \draw[->] (0,-.25)--(0,.3) node[above]{Effect units};
 \draw[->] (0,-.25)--(4.35,-.25) node[right]{Measurement-noise variance};
 \draw[dashed,gray] (0,0)--(4.1,0);
 \draw[thick,ink,dashed] (0,-.2)--(4,-.2);
 \draw[very thick,teal,domain=0:4,samples=60] plot (\x, {-.2+\x/(1+2*\x)});
 \foreach \x in {1,2,3,4}{\draw (\x,-.26)--(\x,-.24);\node[below] at (\x,-.27){\x};}
 \node[left] at (0,-.2){$-0.2$};\node[left] at (0,0){$0$};\node[left] at (0,.2){$0.2$};
 \node[above,align=center] at (2.7,.23){Adjusted regression\ coefficient};
 \node[above] at (2.7,-.2){Specified causal effect};
\end{tikzpicture}
\caption{In the constructed Gaussian example, measurement error can reverse
the sign of an adjusted regression coefficient. The causal effect is fixed at
$\tau=-1/5$; the solid curve is $-1/5+s/(1+2s)$. This is an exact population
calculation, not observed performance of a language model.}
\label{reading:measurement-error}
\end{figure}

The appropriate comparison is consequently broader than extraction accuracy.
Simple dictionaries, supervised extraction, human coding and language models
should use the same permitted document dates and the same causal target.
The evaluation should ask whether their measurement errors change the estimated
effect or the chosen action. It should also test cases in which the notes were
written after treatment or only after a complaint. Those cases alter the
causal role of the text, not merely its predictive quality.
\citep{exec13_census_text_review_arxiv}

\section{Keep the economist's judgment visible}

The literature distinguishes answering questions about a supplied causal
structure from discovering the right structure for a new problem. CLadder
tests formal causal questions under constructed mechanisms. Causal-learn+
concerns assistance with numerical causal analysis. Evidence from either
setting is useful, but it does not measure how reliably an assistant elicits
all relevant variables from an unfamiliar bank process.
\citep{v12_la_2312_04350,v12_la_2606_23608}

The proposal should therefore compare assistance task by task. For formulation,
measure whether it identifies the action, population, timing, outcome and
unresolved assumptions. For diagrams, test whether it preserves ambiguity and
offers competing explanations. For estimation, require the same target and
sampling design as the numerical method. For reporting, ask whether the
reader can distinguish a conditional estimate from an unsupported causal
assertion. The relevant cost includes analyst correction and unresolved cases,
as well as model calls. These are proposed comparisons; the reviewed literature
does not supply a general estimate of savings for this bank.

A successful prototype would make a defined study more reliable or less costly
than credible alternatives. A study in which a simple questionnaire performs
as well would support using the questionnaire. A method that gives fluent
answers while changing the causal question should be rejected for that use.
This decision rule gives language assistance a clear economic purpose without
requiring it to settle causal identification on its own.
